\SourceLeaf{191}
\par\Needspace{6\baselineskip}\medskip
\begin{center}\hypertarget{native-191-article142}{}\pdfbookmark[2]{142.}{article-142}142.\end{center}
Gleichung IV, die in entwickelter Form vom achten Grad wäre, löst man in der vorliegenden Form am schnellsten durch Probieren. Aus der Theorie der Gleichungen lässt sich außerdem leicht zeigen (was wir hier der Kürze halber nicht näher ausführen), dass diese Gleichung entweder zwei oder vier reelle Lösungen zulässt\HistoricalNoteEight{\interlinepenalty=10000 Handschriftliche Bemerkung:] $z$ hat mindestens zwei reelle Werte, da die Werte von $G\sin z^4-\sin(z+H)$ für $z=0$ und für $z=180^0$ entgegengesetzte Vorzeichen haben. $z$ hat nicht mehr als 4 reelle Werte, da $\dfrac{\sin(z+A)}{\sin z^4}$ zwischen $z=0$ und $z=180^0$ nur einmal ein Maximum und einmal ein Minimum erreicht, ebenso zwischen $z=180^0$ und $360^0$. Dies tritt jeweils ein, wenn $2\operatorname{tang}z=-3\operatorname{cotang}A\pm\BareRoot{(9\operatorname{cotang}A^2-16)}=8\operatorname{tang}(z+A)$. —
\par\indent Die Auflösung der Gleichung $a\sin z^4=\sin(z-A)$ lässt sich auf eine zierliche Konstruktion zurückführen. Man setze $\operatorname{cotang}z=x$, also $\sin z=\dfrac{1}{\BareRoot{(1+xx)}}\mathpunct{\cdot}$ Folglich die Gleichung
\[
\frac{1}{(1+xx)^{\frac{3}{2}}}=\frac{\cos A}{a}-\frac{\sin A}{a}x.
\]
Es handelt sich also nun um den Schnitt der Kurve $y=\dfrac{1}{(1+xx)^{\frac{3}{2}}}$ mit der Geraden $y=\dfrac{\cos A}{a}-\dfrac{\sin A}{a}x$.
\par\noindent\textsc{Binet}, \textit{Ecole polyt. Cah. 20. p. 285.}}. Im ersten Fall ist der eine Wert von $\sin z$ positiv, und der negative ist zu verwerfen, da $r'$ nach der Natur der Aufgabe nicht negativ werden kann. Im zweiten Fall ist entweder einer der Werte von $\sin z$ positiv und die übrigen drei sind negativ — dann ist die Wahl eindeutig — oder es sind drei Werte positiv und einer negativ. In diesem letzten Fall sind auch alle etwaigen positiven Werte zu verwerfen, für die $z$ größer als $\delta'$ ist, da nach einer weiteren wesentlichen Bedingung der Aufgabe $\rho'$ und damit auch $\sin(\delta'-z)$ positiv sein muss.

Wenn die Beobachtungen in mäßigen zeitlichen Abständen aufeinanderfolgen, liegt meist der letzte Fall vor, sodass drei positive Werte von $\sin z$ die Gleichung erfüllen. Neben der wahren Lösung findet sich unter ihnen gewöhnlich eine, bei der $z$ nur wenig von $\delta'$ abweicht, bald nach oben, bald nach unten. Diese Erscheinung erklärt sich folgendermaßen. Die analytische Behandlung unserer Aufgabe beruht allein auf der Bedingung, dass die drei Orte des Himmelskörpers im Raum \textit{auf} Geraden liegen müssen, deren Lage durch den absoluten Ort der Erde und die \SourceLeaf{192}beobachtete Position bestimmt ist. Der Natur der Sache nach müssen jene Orte zwar auf \textit{denjenigen} Teilen der Geraden liegen, \textit{von denen} das Licht zur Erde gelangt. Die analytischen Gleichungen kennen diese Einschränkung jedoch nicht. Sie müssen alle Systeme von Orten, die mit den Gesetzen \textsc{Kepler}s vereinbar sind, gleichermaßen erfassen, ob sie nun auf diesen Geraden auf der einen oder anderen Seite der Erde liegen oder mit der Erde selbst zusammenfallen. Dieser letzte Fall genügt unserer Aufgabe gewiss, da sich die Erde selbst nach diesen Gesetzen bewegt. Daraus folgt, dass die Gleichungen eine Lösung umfassen müssen, in der die Punkte $C,C',C''$ mit den Punkten $A,A',A''$ zusammenfallen (sofern man die äußerst kleinen Änderungen der Orte der Erde auf ihrer elliptischen Bahn infolge von Störungen und Parallaxe vernachlässigt). Gleichung IV müsste daher stets die Lösung $z=\delta'$ zulassen, wenn für $P$ und $Q$ die wahren, den Erdorten entsprechenden Werte eingesetzt würden. Setzt man für diese Größen jedoch Werte an, die nur wenig von den wahren Werten abweichen (was bei kurzen Zeitabständen stets angenommen werden darf), so muss unter den Lösungen von Gleichung IV eine liegen, die $z=\delta'$ sehr nahekommt.

Lässt Gleichung IV drei Lösungen mit positiven Werten von $\sin z$ zu, so ergibt die dritte (neben der wahren und der soeben besprochenen) meistens einen Wert von $z$, der größer ist als $\delta'$. Diese Lösung ist daher analytisch möglich, physikalisch jedoch unmöglich. Dann ist die Wahl eindeutig. Die Gleichung kann jedoch zwei verschiedene zulässige Lösungen haben, sodass zwei ganz verschiedene Bahnen unserer Aufgabe genügen. In einem solchen Fall lässt sich die wahre Bahn leicht von der falschen unterscheiden, sobald man weitere, zeitlich weiter entfernte Beobachtungen zur Prüfung heranziehen kann.
